\documentclass[10pt,a4paper]{amsart}
\usepackage[margin=19mm]{geometry}
\usepackage{amsmath,amssymb,mathtools}
\usepackage[scr=rsfs,cal=euler]{mathalfa}
\usepackage{xfrac}
\usepackage[unicode,hidelinks,bookmarks=false]{hyperref}
\newcommand{\ie}{i.e.\ }
\newcommand{\cf}{cf.\ }
\newcommand{\resp}{respectively\ }
\newcommand{\Subsection}{\subsection}
\providecommand{\nobreakdash}{\nobreak\hbox{-}}
\setlength{\emergencystretch}{2em}
\multlinegap0pt
\usepackage{bm}
\usepackage{bbm}
\usepackage{dsfont}
\usepackage{xspace}
\usepackage{cancel}
\usepackage{mathtools}
\usepackage{amsthm}
\makeatletter
\@ifundefined{theorem}{\newtheorem{theorem}{Theorem}}{}
\@ifundefined{lemma}{\newtheorem{lemma}[theorem]{Lemma}}{}
\makeatother
\newcommand{\Aff}{\operatorname{Aff}}
\newcommand{\R}{\mathbb R}
\title[Geometric ideal triangulations of hyperbolic 3-manifolds]{Geometric ideal triangulations of hyperbolic 3-manifolds}
\author{Huabin Ge}
\date{Source: arXiv:2609.27635v1 (CC BY 4.0)}
\begin{document}
\maketitle

\noindent\emph{Source and licence.} Geometric ideal triangulations of hyperbolic 3-manifolds. Version arXiv:2609.27635v1 (CC BY 4.0), distributed under the Creative Commons Attribution 4.0 International licence, as recorded in the arXiv licence metadata for this exact version and shown on the abstract page https://arxiv.org/abs/2609.27635v1. Full rights notice: licenses/arxiv-2609.27635-CC-BY-4.0.txt.\par\smallskip

\noindent\textit{Editorial scope.} Counted unit: the Lemma whose source label is \texttt{lemma:surplus} in the article (source0.tex lines 336--342, source label \texttt{lemma:surplus}), delivered together with the article's own complete original proof of it (source0.tex lines 344--386, one \texttt{proof} environment of 43 lines). No mathematics is written, repaired or re-derived: statement and proof are the article's own text line for line. Required context retained uncounted: none. Every definition, display and label of those lines is retained unchanged. Omitted from this package and not redistributed: 1--335, 343--343, 387--984. Nothing inside a retained excerpt is omitted or altered: every retained body is delivered exactly as the article's lines are written, so no omission receipt of any kind accompanies this package. Citations: the retained excerpts carry no citation command at all, so no bibliography entry of the article is transcribed and no reference is redistributed. Reference adaptation: none was needed and none was made; every reference inside the delivered unit resolves inside it, so no unresolved reference or citation is produced. Printed slips kept literal: no printed word, symbol, hypothesis or number of the counted statement or of its proof was altered. Layout and class: the article's own class is \texttt{article} with packages including geometry, amsmath, amsthm, amsfonts, amssymb, fancyhdr, mathrsfs, graphicx, xy, comment, cases, color, placeins, hyperref; the wrapper is the lane's pinned \texttt{amsart} editorial wrapper at 10pt on A4, which restates the theorem-like environments the retained text needs and adds the article's own macro definitions that the retained text needs (\texttt{\textbackslash Aff}, \texttt{\textbackslash R}), with extra wrapper packages (bm, bbm, dsfont, xspace, cancel, mathtools). The delivered numbering is the unit's own. Assets: no retained span contains a figure environment, a graphics-inclusion command, a PDF-page inclusion, a table or a TikZ picture; no image file is copied into this package, the package declares no asset dependency and no image file is redistributed with it. Licence: version arXiv:2609.27635v1 of this article is distributed under the Creative Commons Attribution 4.0 International licence, as recorded in the arXiv licence metadata for this exact version and shown on the abstract page https://arxiv.org/abs/2609.27635v1. A full rights notice is delivered at licenses/arxiv-2609.27635-CC-BY-4.0.txt. Characters: every retained excerpt is pure ASCII, so the delivered engine, XeLaTeX with its default Latin Modern text font, reports no missing character.
\begin{lemma}\label{lemma:surplus}
Let $f_i$ be the number of two-dimensional faces of $P_i$. The vector space of compatible height classes on $\mathcal P$ has dimension at least
\begin{equation}\label{equation:surplus}
\frac12\sum_i(f_i-4).
\end{equation}
In particular, if one of the $P_i$ is not a tetrahedron, there is a nonzero compatible height class.
\end{lemma}

\begin{proof}
Let $V_i=V(P_i)$ and set
\begin{equation}
\Aff(P_i)=\{r|_{V_i}:r:A_i\to\R\text{ affine}\},
\qquad
H_i=\R^{V_i}/\Aff(P_i).
  \nonumber \end{equation}
Since $A_i$ has dimension three,
\begin{equation}
\dim H_i=|V_i|-4.
  \nonumber \end{equation}
For a paired face $F^-_\alpha$ with $n_\alpha$ vertices, let
\begin{equation}
\Aff(F^-_\alpha)=\{r|_{V(F^-_\alpha)}:r:\operatorname{aff}(F^-_\alpha)\to\R\text{ affine}\}
  \nonumber \end{equation}
and put
\begin{equation}
Q_\alpha=\R^{V(F^-_\alpha)}/\Aff(F^-_\alpha),
\qquad
\dim Q_\alpha=n_\alpha-3.
  \nonumber \end{equation}
The face pairings define a linear map
\begin{equation}
B:\bigoplus_i H_i\longrightarrow\bigoplus_\alpha Q_\alpha,
\qquad
B_\alpha([h])=[h^-_\alpha-h^+_\alpha\circ\phi_\alpha].
  \nonumber \end{equation}
It is well defined because the restriction of an affine function remains affine after composition with an affine homeomorphism. The compatible height classes are precisely $\ker B$.

Write $v_i,e_i,f_i$ for the numbers of vertices, edges and two-dimensional faces of $P_i$, and let $F$ be the number of pairs of faces. Then
\begin{equation}
\sum_i f_i=2F,
\qquad
\sum_i e_i=\sum_\alpha n_\alpha.
  \nonumber \end{equation}
Euler's formula gives
\begin{equation}
\dim\bigoplus_iH_i-\dim\bigoplus_\alpha Q_\alpha
 =\sum_i(v_i-4)-\sum_\alpha(n_\alpha-3)
 =\frac12\sum_i(f_i-4).
  \nonumber \end{equation}
Hence $\dim\ker B$ is at least the number in \eqref{equation:surplus}. Since a three-dimensional convex polyhedron has at least four faces, with equality only for a tetrahedron, the last statement follows.
\end{proof}

\end{document}
